\textbf{1. Kiến thức cần nhớ}
\begin{itemize}
\item \textbf{Số gần đúng:} Trong đo đạc và tính toán, ta thường nhận được số gần đúng $a$ của số đúng $\overline{a}$.
\item \textbf{Sai số tuyệt đối} của số gần đúng $a$: $\boxed{\Delta_a = |\overline{a} - a|}$.
\item \textbf{Độ chính xác:} Nếu $\Delta_a = |\overline{a} - a| \le d$ thì $a - d \le \overline{a} \le a + d$. Ta nói $a$ là số gần đúng của $\overline{a}$ với độ chính xác $d$, viết gọn là $\boxed{\overline{a} = a \pm d}$.
\item \textbf{Sai số tương đối} của số gần đúng $a$: $\boxed{\delta_a = \dfrac{\Delta_a}{|a|}}$. Nếu $\overline{a} = a \pm d$ thì $\delta_a \le \dfrac{d}{|a|}$. Sai số tương đối càng nhỏ thì chất lượng phép đo đạc, tính toán càng cao (thường viết dưới dạng phần trăm).
\item \textbf{Quy tắc làm tròn số:} Xét chữ số ngay sau hàng quy tròn: nếu nhỏ hơn $5$ thì thay nó và các chữ số bên phải nó bởi chữ số $0$; nếu lớn hơn hoặc bằng $5$ thì cộng thêm $1$ vào chữ số hàng quy tròn rồi thay các chữ số bên phải nó bởi chữ số $0$ (bỏ các chữ số $0$ ở tận cùng bên phải phần thập phân).
\item \textbf{Quy tròn số gần đúng với độ chính xác $d$:} Quy tròn số $a$ đến hàng lớn hơn một bậc so với hàng của chữ số khác $0$ đầu tiên bên trái của $d$.
\end{itemize}

\textbf{2. Phương pháp giải các dạng}

%% DANG 0D6?1-1
\textbf{Dạng 1. Tính và ước lượng sai số tuyệt đối, tương đối}
\begin{enumerate}
\item Xác định số đúng $\overline{a}$, số gần đúng $a$ (hoặc chặn khoảng $m \le \overline{a} \le M$).
\item Tính hoặc ước lượng sai số tuyệt đối $\Delta_a = |\overline{a} - a|$ (nếu biết $m \le \overline{a} \le M$ thì $\Delta_a \le \max\{|a-m|; |M-a|\}$).
\item Tính hoặc ước lượng sai số tương đối $\delta_a = \dfrac{\Delta_a}{|a|} \le \dfrac{d}{|a|}$ và đổi sang phần trăm nếu đề yêu cầu.
\end{enumerate}
Lưu ý: Để so sánh độ chính xác của hai phép đo có đơn vị hoặc độ lớn khác nhau, phải so sánh sai số tương đối $\delta_a$, không so sánh sai số tuyệt đối $\Delta_a$.

%% DANG 0D6?1-2
\textbf{Dạng 2. Tính và xác định độ chính xác của kết quả}
\begin{enumerate}
\item Từ giả thiết, viết kết quả đo hoặc tính toán dưới dạng $m \le \overline{a} \le M$ hoặc tìm cận trên của $|\overline{a} - a|$.
\item Nếu biết $a$ là số gần đúng, độ chính xác là số $d > 0$ thỏa mãn $|\overline{a} - a| \le d$.
\item Nếu chỉ biết $m \le \overline{a} \le M$, chọn $a = \dfrac{m+M}{2}$ và độ chính xác $d = \dfrac{M-m}{2}$ để viết $\overline{a} = a \pm d$.
\end{enumerate}
Lưu ý: Số $d$ càng nhỏ thì độ chính xác càng cao; mọi số $d' > d$ cũng là một độ chính xác của $a$.

%% DANG 0D6?1-3
\textbf{Dạng 3. Quy tròn số gần đúng}
\begin{enumerate}
\item Xác định hàng quy tròn: nếu đề cho sẵn hàng thì dùng ngay; nếu đề cho độ chính xác $d$ thì tìm hàng của chữ số khác $0$ đầu tiên bên trái của $d$, rồi chọn hàng liền trước nó (gấp $10$ lần hàng đó).
\item Áp dụng quy tắc làm tròn số tại hàng quy tròn vừa tìm được.
\end{enumerate}
Lưu ý: Dễ nhầm lẫn giữa hàng của chữ số khác $0$ đầu tiên của $d$ với hàng quy tròn (ví dụ: $d = 0{,}002$ ở hàng phần nghìn thì phải quy tròn đến hàng phần trăm).

%% DANG 0D6?1-4
\textbf{Dạng 4. Viết số gần đúng cho số đúng biết độ chính xác}
\begin{enumerate}
\item Dùng máy tính cầm tay (hoặc tính toán) tìm giá trị thập phân của số đúng $\overline{a}$ với đủ chữ số thập phân.
\item Xác định hàng của chữ số khác $0$ đầu tiên bên trái của độ chính xác $d$.
\item Làm tròn số $\overline{a}$ đến hàng lớn hơn một bậc (hoặc đến hàng sao cho sai số làm tròn không vượt quá $d$) để thu được số gần đúng $a$.
\end{enumerate}
Lưu ý: Không bỏ quên các chữ số $0$ ở phần nguyên khi làm tròn đến hàng chục, hàng trăm, hàng nghìn (ví dụ: $12\,480$ làm tròn đến hàng trăm là $12\,500$, không phải $125$).
